%% ma: L12-B08-0096
%% lop: 12
%% chuong: 2
%% bai: 8
%% dang: 12.8.29
%% loai: TN4
%% muc: NB
%% dap_an: "B"
%% nguon: "CĐ2.3.pdf giáo viên gửi 27/10/2026 (M07, Lê Minh Khôi – Nguyễn Anh Khôi, Chương II Chủ đề 3 Biểu thức tọa độ của các phép toán vectơ), Câu 4"
%% kien_thuc: "Bài 8 (biểu thức tọa độ của phép cộng, trừ vectơ, tích vectơ với một số)"
%% trang_thai: chinh-thuc
\begin{ex}
Trong không gian $Oxyz$, cho hai vectơ $\vec a=(0;7;2)$ và $\vec b=(-3;4;1)$. Tọa độ của vectơ $\vec p=2\cdot\vec a-3\cdot\vec b$ là
\choice{$\vec p=(-9;26;7)$}{\True $\vec p=(9;2;1)$}{$\vec p=(3;3;1)$}{$\vec p=(3;10;3)$}
\loigiai{$\vec p=(0+9;14-12;4-3)=(9;2;1)$.}
\end{ex}
